\documentclass[10pt,a4paper]{amsart}
\usepackage[margin=19mm]{geometry}
\usepackage{amsmath,amssymb,mathtools}
\usepackage[scr=rsfs,cal=euler]{mathalfa}
\usepackage{xfrac}
\usepackage[unicode,hidelinks,bookmarks=false]{hyperref}
\newcommand{\ie}{i.e.\ }
\newcommand{\cf}{cf.\ }
\newcommand{\resp}{respectively\ }
\newcommand{\Subsection}{\subsection}
\providecommand{\nobreakdash}{\nobreak\hbox{-}}
\setlength{\emergencystretch}{2em}
\multlinegap0pt
\usepackage{amssymb}
\newtheorem{theorem}{Theorem}
\title{On $\phi$-$(k,n)$-absorbing $\delta$-primary hyperideals}
\author{Mahdi Anbarloei}
\date{Source: arXiv:2607.27223v1 (CC BY 4.0)}
\begin{document}
\maketitle

\noindent\emph{Source and licence.} On $\phi$-$(k,n)$-absorbing $\delta$-primary hyperideals. Version arXiv:2607.27223v1 (CC BY 4.0), distributed by arXiv under the Creative Commons Attribution 4.0 International licence, http://creativecommons.org/licenses/by/4.0/, as recorded in the arXiv licence metadata for this exact version and shown on the abstract page https://arxiv.org/abs/2607.27223v1. Full licence notice: \texttt{licenses/arxiv-2607.27223-CC-BY-4.0.txt}.\par\smallskip

\noindent\textit{Editorial scope.} Counted unit: the article's theorem weak (source0.tex lines 282--284). The article's complete original proof of it is delivered with it (source0.tex lines 285--351, one \texttt{proof} environment of 67 lines). No mathematics is written, repaired or re-derived: the statement and the proof are the article's own lines, in the article's own notation, and no retained line is substituted or re-flowed.  No further source-local context is needed: every cross-reference of every retained line resolves inside the two delivered spans.  Nothing was omitted from any retained span, so the package carries no omission record.  No retained line contains a citation, so the wrapper carries no bibliography and no bibliography entry.  No retained line contains a figure environment, an image-inclusion command or drawing code, so no image is delivered and the package declares no asset dependency.  Editorial formatting only, in addition: an AMS \texttt{amsart} preamble that restates the article's own declarations and macros used by the retained lines, namely the environments \texttt{theorem} on separate counters, together with the standard packages those lines need.  The retained environments print in this document as Theorem 1 in order of appearance, whereas the article prints the same items with the numbering of its own full document.  The complete native source of the article as distributed in the arXiv e-print for this exact version is bundled unchanged as \texttt{sources/206339/source0.tex}.
\begin{theorem} \label{weak}
Let $I$ be a proper hyperideal of $G$. Then $I$  is   a  $\phi$-$(2,2)$-absorbing $\delta$-primary hyperideal of $G$ if and only if   $I/\phi(I)$ is a  weakly $(2,2)$-absorbing $\delta_q$-primary hyperideal of $G/\phi(I)$. 
\end{theorem}

\begin{proof}
 $\Longrightarrow$ Let  $I$  be  a strongly $\phi$-$(2,2)$-absorbing $\delta$-primary hyperideal of $G$ and \\

$\phi(I) \neq g(f(t_{11}^{1(i-1)},\phi(I),t_{1(i+1)}^{1m}),f(t_{21}^{2(i-1)},\phi(I),t_{2(i+1)}^{2m}),f(t_{31}^{3(i-1)},\phi(I),t_{3(i+1)}^{3m}))$

$\hspace{0.7cm}=f(g(t_{11}^{31}),\ldots,g(t_{1(i-1)}^{3(i-1)}),\phi(I),g(t_{1(i+1)}^{3(i+1)}),\ldots,g(t_{1m}^{3m}))$

$\hspace{0.7cm} \in I/\phi(I)$. 

for $t_{11}^{1m},t_{21}^{2m},t_{31}^{3m}\in G$.
Then we conclude that

$\hspace{1cm} f(g(t_{11}^{31}),\ldots, g(t_{1(i-1)}^{3(i-1)}),0,g(t_{1(i+1)}^{3(i+1)}),\ldots,g(t_{1m}^{3m}))$  

$\hspace{1cm}=g(f(t_{11}^{1(i-1)},0,t_{1(i+1)}^{1m}),f(t_{21}^{2(i-1)},0,t_{2(i+1)}^{2m}),f(t_{31}^{3(i-1)},0,t_{3(i+1)}^{3m}))$

$\hspace{1cm}\subseteq  I-\phi(I).$\\
This implies that

 $\hspace{1cm}g(f(t_{11}^{1(i-1)},0,t_{1(i+1)}^{1m}),f(t_{21}^{2(i-1)},0,t_{2(i+1)}^{2m}))$

$ \hspace{1.5cm}=f(g(t_{11}^{21}),\ldots,g(t_{1(i-1)}^{2(i-1)}),0,g(t_{1(i+1)}^{2(i+1)}),\ldots,g(t_{1m}^{2m}))\subseteq I$\\ or 

$\hspace{1cm}g(f(t_{21}^{2(i-1)},0,t_{2(i+1)}^{2m}),f(t_{31}^{3(i-1)},0,t_{3(i+1)}^{3m}))$

 $\hspace{1.5cm}=f(g(t_{21}^{31}),\ldots,g(t_{2(i-1)}^{3(i-1)}),0,g(t_{2(i+1)}^{3(i+1)}),\ldots,g(t_{2m}^{3m}))\subseteq \delta(I)$\\
or 


$\hspace{1cm}g(f(t_{11}^{1(i-1)},0,t_{1(i+1)}^{1m}),f(t_{31}^{3(i-1)},0,t_{3(i+1)}^{3m}))$

$\hspace{1.5cm}=f(g(t_{11}^{31}),\ldots,g(t_{1(i-1)}^{3(i-1)}),0,g(t_{1(i+1)}^{3(i+1)}),\ldots,g(t_{1m}^{3m}))\subseteq \delta(I)$.\\ 
as $I$ is a strongly $\phi$-$(2,2)$-absorbing $\delta$-primary hyperideal of $G$. Therefore, we get the result that 

$\hspace{1cm}f(g(t_{11}^{21}),\ldots,g(t_{1(i-1)}^{2(i-1)}),\phi(I),g(t_{1(i+1)}^{2(i+1)}),\ldots,g(t_{1m}^{2m}))$


$\hspace{1.5cm}=g(f(t_{11}^{1(i-1)},\phi(I),t_{1(i+1)}^{1m}),f(t_{21}^{2(i-1)},\phi(I),t_{2(i+1)}^{2m}))\in I/\phi(I)$\\ or  

$\hspace{1cm}f(g(t_{21}^{31}),\ldots,g(t_{2(i-1)}^{3(i-1)}),\phi(I),g(t_{2(i+1)}^{3(i+1)}),\ldots,g(t_{2m}^{3m}))$

$\hspace{1.5cm}=g(f(t_{21}^{2(i-1)},\phi(I),t_{2(i+1)}^{2m}),f(t_{31}^{3(i-1)},\phi(I),t_{3(i+1)}^{3m}))$

$\hspace{1.5cm}\in \delta(I)/\phi(I)=\delta_q(I/\phi(I))$\\
or 

$\hspace{1cm}f(g(t_{11}^{31}),\ldots,g(t_{1(i-1)}^{3(i-1)}),\phi(I),g(t_{1(i+1)}^{3(i+1)}),\ldots,g(t_{1m}^{3m}))$

$\hspace{1.5cm}=g(f(t_{11}^{1(i-1)},\phi(I),t_{1(i+1)}^{1m}),f(t_{31}^{3(i-1)},\phi(I),t_{3(i+1)}^{3m}))$

$ \hspace{1.5cm} \in \delta(I)/\phi(I)=\delta_q(I/\phi(I))$.\\
Consequently, $I/\phi(I)$ is a  weakly $(2,2)$-absorbing $\delta_q$-primary hyperideal of $G/\phi(I)$. 

$\Longleftarrow $ Assume that  $g(t_1^3) \in I-\phi(I)$ for  $t_1^3 \in G$. Then $f(g(t_1^3),\phi(I),0^{(m-2)}) \neq \phi(I)$. Therefore we conclude that

$\phi(I) \neq g(f(t_1,\phi(I),0^{(m-2)}),f(t_2,\phi(I),0^{(m-2)}),f(t_3,\phi(I),0^{(m-2)})) \in I/\phi(I)$. \\
 Since $I/\phi(I)$ is a  weakly $(2,2)$-absorbing $\delta_q$-primary hyperideal of $G/\phi(I)$, we have 
 
 $g(f(t_1,\phi(I),0^{(m-2)}),f(t_2,\phi(I),0^{(m-2)}))=f(g(t_1^2),\phi(I),0^{(m-2)}) \in I/\phi(I)$.\\
or 
 
 $g(f(t_2,\phi(I),0^{(m-2)}),f(t_3,\phi(I),0^{(m-2)}))  =f(g(t_2^3),\phi(I),0^{(m-2)}) \in \delta(I)/\phi(I)$.\\
or 
 
$g(f(t_1,\phi(I),0^{(m-2)}),f(t_3,\phi(I),0^{(m-2)})) =f(g(t_1^3),\phi(I),0^{(m-2)})  \in \delta(I)/\phi(I)$.\\
Hence we obtain  $g(t_1^2) \in I$ or $g(t_2^3) \in \delta(I)$ or $g(t_1^3) \in \delta(I)$. Thus, $I$  is   a $\phi$-$(2,2)$-absorbing primary hyperideal of $G$. 
\end{proof}

\end{document}
